\section{Missing Proofs for Known Implications}
\label{sec:known-impls-extra}

\thmRecall{rthm:impl:ext-to-epistemic}
\begin{proof}
Suppose $\phi(F_1, F_2)$.
Pick any $(\Ical, A) \in \Omega$.
Let $\Ical \defeq ([n], [m], (v_i)_{i=1}^n, w)$.
Pick any $i \in [n]$.

Suppose $A$ is epistemic-$F_1$-fair to $i$.
Let $B$ be $i$'s epistemic-$F_1$-certificate for $A$.
Then $B$ is $F_1$-fair to $i$ and $A_i = B_i$.
By $\phi(F_1, F_2)$, $B$ is also $F_2$-fair to $i$.
Hence, $B$ is $i$'s epistemic-$F_2$-certificate for $A$.
Therefore, $\phi(\textrm{epistemic-}F_1, \textrm{epistemic-}F_2)$ holds.

Suppose $A$ is min-$F_1$-share-fair to $i$.
Let $B$ be $i$'s min-$F_1$-share-certificate for $A$.
Then $B$ is $F_1$-fair to $i$ and $v_i(A_i) \ge v_i(B_i)$.
By $\phi(F_1, F_2)$, $B$ is also $F_2$-fair to $i$.
Hence, $B$ is $i$'s min-$F_2$-share-certificate for $A$.
Therefore, $\phi(\textrm{min-}F_1\textrm{-share}, \textrm{min-}F_2\textrm{-share})$ holds.
\end{proof}

\thmRecall{rthm:impl:propx-to-propm}
\begin{proof}
Assume (for the sake of contradiction) that there is an allocation $A$ where
agent $i$ is PROPx-satisfied but not PROPm-satisfied.
%
Since $i$ is not PROPm-satisfied, we get $v_i(A_i) \le w_iv_i([m])$.
Since $i$ is PROPx-satisfied, we get

\begin{itemize}
\item $v_i(A_i \setminus S) > w_iv_i([m])$ for all $S \subseteq A_i$ such that
    $v_i(S \mid A_i \setminus S) < 0$.
\item $v_i(A_i \cup S) > w_iv_i([m])$ for all $S \subseteq [m] \setminus A_i$
    such that $v_i(S \mid A_i) > 0$.
\end{itemize}

Since $i$ is not PROPm-satisfied, we get that $T \neq \emptyset$ and $v_i(A_i) + \max(T) \le w_iv_i([m])$.
Let $\max(T) = \tau_{\jhat} = v_i(\Shat \mid A_i) > 0$.
Then $v_i(A_i \cup \Shat) \le w_iv_i([m])$, which contradicts the fact that
$i$ is PROPx-satisfied by $A$.
Hence, if $i$ is PROPx-satisfied by $A$, then she is also PROPm-satisfied by $A$.
\end{proof}

\thmRecall{rthm:impl:mefs-to-prop}
\begin{proof}
Let $B$ be agent $i$'s MEFS-certificate for $A$.
Then for all $j \in [n]$, we have $v_i(B_i)/w_i \ge v_i(B_j)/w_j$.
Sum these inequalities over all $j \in [n]$, weighting each by $w_j$,
to get $v_i(B_i)/w_i \ge \sum_{j=1}^n v_i(B_j)$.
Since $v_i$ is subadditive, we get $v_i([m]) \le \sum_{j=1}^n v_i(B_j)$.
Hence,
\[ \frac{v_i(A_i)}{w_i} \ge \frac{v_i(B_i)}{w_i} \ge \sum_{j=1}^n v_i(B_j) \ge v_i([m]). \]
\end{proof}

\thmRecall{rthm:impl:ef-to-gprop}
\begin{proof}
Let $S$ be a subset of agents containing $i$.
Let $\Ahat$ be the allocation obtained by restricting $A$ to $S$ (c.f.~\cref{defn:restricting}).
Then $i$ is also envy-free in $\Ahat$.
$i$ is also MEFS-satisfied by $\Ahat$, since $\Ahat$ is its own MEFS-certificate for agent $i$.
By \cref{thm:impl:mefs-to-prop}, agent $i$ is PROP-satisfied by $\Ahat$.
Since this is true for all $S$ containing $i$,
we get that $A$ is groupwise-PROP-fair to agent $i$.
\end{proof}

\thmRecall{rthm:impl:prop-to-ef-superadd-id}
\begin{proof}
Let $v$ be the common valuation function. Let $A$ be a PROP allocation.
Then for each agent $i$, we have $v(A_i) \ge w_iv([m])$.
Suppose $v(A_k) > w_kv([m])$ for some agent $k$.
Sum these inequalities to get $\sum_{i=1}^n v(A_i) > v([m])$.
This contradicts superadditivity of $v$, so $v(A_i) = w_iv([m])$ for each agent $i$.
Hence, $v(A_i)/w_i = v([m])$ for all $i$, so $A$ is EF.
\end{proof}

\thmRecall{rthm:impl:prop-to-ef-n2}
\begin{proof}
Assume $i=1$ \wLoG. Then
\begin{align*}
& \frac{v_1(A_1)}{w_1} \ge v_1([m]) \ge v_1(A_1) + v_1(A_2)
\\ &\implies v_1(A_2) \le v_1(A_1)\left(\frac{1}{w_1} - 1\right) = w_2 \frac{v_1(A_1)}{w_1}.
\end{align*}
Hence, agent 1 does not envy agent 2.
\end{proof}

\thmRecall{rthm:impl:efx-to-ef1}
\begin{proof}
Suppose $i$ is EFX-satisfied but not EF1-satisfied.
Suppose $i$ EF1-envies $j$.

Since $i$ EF1-envies $j$, we get that for all $t \in A_j$, we have
\[ \frac{v_i(A_i)}{w_i} < \frac{v_i(A_j \setminus \{t\})}{w_j}. \]
Since $i$ is EFX-saisfied, we get that
for all $t \in A_j$ such that $v_i(t \mid A_i) > 0$, we have
\[ \frac{v_i(A_i)}{w_i} \ge \frac{v_i(A_j \setminus \{t\})}{w_j}. \]
Hence, for all $t \in A_j$, we get $v_i(t \mid A_i) \le 0$.

Since $i$ EF1-envies $j$, we get that for all $t \in A_i$, we have
\[ \frac{v_i(A_i \setminus \{t\})}{w_i} < \frac{v_i(A_j)}{w_j}. \]
Since $i$ is EFX-saisfied, we get that
for all $t \in A_i$ such that $v_i(t \mid A_i \setminus \{t\}) < 0$, we have
\[ \frac{v_i(A_i \setminus \{t\})}{w_i} \ge \frac{v_i(A_j)}{w_j}. \]
Hence, for all $t \in A_i$, we get $v_i(t \mid A_i \setminus \{t\}) \ge 0$.

If $v_i$ is additive, we get $v_i(A_i) \ge 0 \ge v_i(A_j)$, which is a contradiction.

If $v_i$ is doubly-monotone and $v_i(g \mid S) > 0$ for every good $g$
and $v_i(c \mid S) < 0$ for every chore $c$, then
all items in $A_j$ are chores and all items in $A_i$ are goods.
Hence, $v_i(A_i) \ge 0 \ge v_i(A_j)$, which is a contradiction.

Suppose all agents have equal entitlements, all items are goods for $i$, and $v_i$ is submodular.
Let $A_j \defeq \{g_1, \ldots, g_k\}$. Then
\[ v_i(A_j \mid A_i) = \sum_{t=1}^k v_i(g_t \mid A_i \cup \{g_1, \ldots, g_{t-1}\})
\le \sum_{t=1}^k v_i(g_t \mid A_i) \le 0. \]
Hence, $v_i(A_j) \le v_i(A_i \cup A_j) = v_i(A_i) + v_i(A_j \mid A_i) \le v_i(A_i)$.
This is a contradiction.

Suppose all items are chores for $i$ and $v_i$ is submodular.
Let $A_i = \{c_1, \ldots, c_k\}$. Then
\[ v_i(A_i) = \sum_{t=1}^k v_i(c_t \mid \{c_1, \ldots, c_{t-1}\})
\ge \sum_{t=1}^k v_i(c_t \mid A_i \setminus \{c_t\}) \ge 0. \]
Hence, $v_i(A_i) \ge 0 \ge v_i(A_j)$, which is a contradiction.

A contradiction implies that it's impossible for agent $i$ to be
EFX-satisfied but not EF1-satisfied.
\end{proof}

\thmRecall{rthm:mms-and-all-envy}
\begin{proof}
Suppose agent $i$ is not EFX-satisfied by $A$, i.e., she EFX-envies some agent $j$.
Then $\exists S \subseteq A_i \cup A_j$ where either

\begin{tightenum}
\item $S \subseteq A_j$, $v_i(S \mid A_i) > 0$, and
    $\displaystyle \frac{v_i(A_i)}{w_i} < \frac{v_i(A_j \setminus S)}{w_j}$.
\item $S \subseteq A_i$, $v_i(S \mid A_i \setminus S) < 0$, and
    $\displaystyle \frac{v_i(A_i \setminus S)}{w_i} < \frac{v_i(A_j)}{w_j}$.
\end{tightenum}

If all items are goods, case 2 doesn't occur.

\textbf{Case 1}: $S \subseteq A_j$

Let $B$ be the allocation obtained by transferring $S$ from $A_j$ to $A_i$.
Formally, let $B_i \defeq A_i \cup S$, $B_j \defeq A_j \setminus S$,
and $B_k \defeq A_k$ for all $k \in [n] \setminus \{i, j\}$. Then
\[ \frac{v_i(B_i)}{w_i} = \frac{v_i(A_i) + v_i(S \mid A_i)}{w_i} > \frac{v_i(A_i)}{w_i}, \]
\[ \frac{v_i(B_j)}{w_j} = \frac{v_i(A_j \setminus S)}{w_j} > \frac{v_i(A_i)}{w_i}, \]
and for any $k \in [n] \setminus \{i, j\}$, we get
\[ \frac{v_i(B_k)}{w_k} = \frac{v_i(A_k)}{w_k} > \frac{v_i(A_i)}{w_i}. \]
Hence,
\[ \min_{k=1}^n \frac{v_i(B_k)}{w_k} > \frac{v_i(A_i)}{w_i} \ge \frac{\WMMS_i}{w_i}, \]
which is a contradiction.

\textbf{Case 2}: $S \subseteq A_i$

Let $v_i$ be additive and $w_i \le w_j$.
Let $B$ be the allocation obtained by transferring $S$ from $A_i$ to $A_j$. Formally,
let $B_i \defeq A_i \setminus S$, $B_j \defeq A_j \cup S$,
and $B_k \defeq A_k$ for all $k \in [n] \setminus \{i, j\}$. Then
\[ \frac{v_i(B_i)}{w_i} = \frac{v_i(A_i) - v_i(S)}{w_i} > \frac{v_i(A_i)}{w_i}, \]
\[ \frac{v_i(B_j)}{w_j} = \frac{v_i(A_j) + v_i(S)}{w_j} > \frac{v_i(A_i \setminus S)}{w_i} - \frac{d_i(S)}{w_j} \ge \frac{v_i(A_i)}{w_i}, \]
and for any $k \in [n] \setminus \{i, j\}$, we get
\[ \frac{v_i(B_k)}{w_k} = \frac{v_i(A_k)}{w_k} > \frac{v_i(A_i)}{w_i}. \]
Hence,
\[ \min_{k=1}^n \frac{v_i(B_k)}{w_k} > \frac{v_i(A_i)}{w_i} \ge \frac{\WMMS_i}{w_i}, \]
which is a contradiction.
\end{proof}

\thmRecall{rthm:impl:mms-to-eefx}
\begin{proof}
In any allocation $X$, let
$E_X$ be the set of agents envied by $i$,
$S_X$ be the set of agents EFX-envied by $i$,
$W_X$ be the total number of items among the agents in $S_X$.

$E_X \defeq \{t \in [n] \setminus \{i\}: i \textrm{ envies } t \textrm{ in } X\}$,
$S_X \defeq \{t \in [n] \setminus \{i\}: i \textrm{ EFX-envies } t \textrm{ in } X\}$,
$W_X \defeq \sum_{t \in S_X} |X_t|$, and $\phi(X) \defeq (-|E_X|, W_X)$.

First, we show that for any allocation $X$ where $|E_X| \le n-2$ and $S_X \neq \emptyset$,
there exists a \emph{better} allocation $Y$, i.e, $Y_i = X_i$ and $\phi(Y) < \phi(X)$
(tuples are compared lexicographically).

Let $j \in S_X$ and $k \in [n] \setminus \{i\} \setminus E_X$.
Since $i$ EFX-envies $j$, $\exists S \subseteq X_j$ such that $v_i(S \mid X_i) > 0$ and
\[ \frac{v_1(X_1)}{w_1} < \frac{v_1(X_j \setminus S)}{w_j}. \]
Let $Y_k \defeq X_k \cup S$, $Y_j \defeq X_j \setminus S$,
and $Y_t \defeq X_t$ for all $t \in [n] \setminus \{j, k\}$.

Then $i$ envies $j$ in $Y$. Hence, $E_X \subseteq E_Y$.
If $k \in E_Y$, then $|E_Y| > |E_X|$, so $\phi(Y) < \phi(X)$.
If $k \not\in E_Y$, then $W_Y < W_X$, so $\phi(Y) < \phi(X)$.
Hence, $Y$ is better than $X$.

Set $X = A$. As long as $|E_X| < n-1$ and $S_X \neq \emptyset$,
keep modifying $X$ as per Lemma 1.
This process will eventually end, since $\phi(X)$ keeps reducing,
and there are a finite number of different values $\phi(X)$ can take.
Let $B$ be the final allocation thus obtained.
Then $B_i = A_i$, and $|E_B| = n-1$ or $S_B = \emptyset$.

By <a href="../misc/mms-and-all-envy.html">mms-and-all-envy.html</a>,
$|E_B| = n-1$ implies $S_B = \emptyset$.
Hence, $B$ is agent $i$'s EEFX-certificate for $A$.
\end{proof}

\thmRecall{rthm:submod-marginal-is-submod}
\begin{proof}
Let $P, Q \subseteq \Omega \setminus X$. Let $g(Y) \defeq f(Y \mid X)$. Then
\begin{align*}
& g(P) + g(Q)
\\ &= f(P \cup X) + f(Q \cup X) - 2f(X)
\\ &\ge f((P \cup X) \cup (Q \cup X)) + f((P \cup X) \cap (Q \cap X)) - 2f(X)
    \tag{by $f$'s submodularity}
\\ &= f((P \cup Q) \cup X) + f((P \cap Q) \cup X) - 2f(X)
    \tag{by De Morgan's law}
\\ &= g(P \cup Q) + g(P \cap Q).
\end{align*}
Hence, $g$ is submodular.
\end{proof}

\thmRecall{rthm:impl:propm-to-prop1}
\begin{proof}
Suppose allocation $A$ is PROPm-fair to $i$ but not PROP1-fair to $i$. Then
\begin{tightenum}
\item\label{item:impl:propm-to-prop1:unprop}$v_i(A_i) < w_iv_i([m])$ (by PROP1 unfairness).
\item\label{item:impl:propm-to-prop1:unprop1-chores}$v_i(A_i \setminus \{c\}) \le w_iv_i([m])$
    for all $c \in A_i$ (by PROP1 unfairness).
\item\label{item:impl:propm-to-prop1:unprop1-goods}$v_i(A_i \cup \{g\}) \le w_iv_i([m])$
    for all $g \in [m] \setminus A_i$ (by PROP1 unfairness).
\item\label{item:impl:propm-to-prop1:propm-chores}$v_i(A_i \setminus \{c\}) > w_iv_i([m])$
    for all $c \in A_i$ such that $v_i(c \mid A_i \setminus \{c\}) < 0$ (by PROPm fairness).
\item\label{item:impl:propm-to-prop1:propm-goods}$T = \emptyset$ or $v_i(A_i) + \max(T) > w_iv_i([m])$
    (by PROPm fairness; c.f.~\cref{defn:propm} for the definition of $T$).
\end{tightenum}

By \ref{item:impl:propm-to-prop1:unprop1-chores} and \ref{item:impl:propm-to-prop1:propm-chores},
we get $v_i(c \mid A_i \setminus \{c\}) \ge 0$ for all $c \in A_i$.
We now show that $v_i(A_i) \ge 0$.
If $v_i$ is doubly strictly monotone, then $A_i$ only contains goods, so $v_i(A_i) \ge 0$.
Now suppose $v_i$ is submodular. Let $A_i = \{g_1, \ldots, g_k\}$. Then
\[ v_i(A_i) = \sum_{j=1}^k v_i(g_j \mid \{g_1, \ldots, g_{j-1}\})
    \ge v_i(g_j \mid A_i \setminus \{g_j\}) \ge 0. \]

Suppose $T = \emptyset$. Then $\tau_j = 0$ for all $j \in [n] \setminus \{i\}$.
Hence, for all $j \in [n] \setminus \{i\}$, we have $v_i(A_j \mid A_i) \le 0$.
If $v_i$ is doubly strictly monotone, then $[m] \setminus A_i$ contains only chores,
so $v_i([m] \setminus A_i \mid A_i) \le 0$. If $v_i$ is submodular, then
$v_i(\cdot \mid A_i)$ is subadditive by \cref{thm:submod-marginal-is-submod}, so
\[ v_i([m] \setminus A_i \mid A_i)
    \le \sum_{j \in [n] \setminus \{i\}} v_i(A_j \mid A_i) \le 0. \]
Hence, $v_i(A_i) \ge v_i([m])$.
If $v_i([m]) \le 0$, then $v_i(A_i) \ge 0 \ge w_iv_i([m])$,
and if $v_i([m]) \ge 0$, then $v_i(A_i) \ge v_i([m]) \ge w_iv_i([m])$.
This contradicts \ref{item:impl:propm-to-prop1:unprop}, so $T \neq \emptyset$.

Let $\max(T) = \tau_{\jhat} > 0$. By definition of $\tau_{\jhat}$, we get
\begin{align*}
0 < \tau_{\jhat} &= \min(\{v_i(S \mid A_i) \mid S \subseteq A_{\jhat} \textrm{ and } v_i(S \mid A_i) > 0\})
\\ &\le \min(\{v_i(g \mid A_i) \mid g \in A_{\jhat} \textrm{ and } v_i(g \mid A_i) > 0\}).
\end{align*}

\textbf{Case 1}: $v_i(g \mid A_i) \le 0$ for all $g \in A_{\jhat}$.
\\ If $v_i$ is doubly strictly monotone, then $A_{\jhat}$ only has chores,
and so $v_i(S \mid A_i) \le 0$ for all $S \subseteq A_{\jhat}$.
This contradicts the fact that $\tau_{\jhat} > 0$.
Now let $v_i$ be submodular.
Since $v_i(\cdot \mid A_i)$ is subadditive by \cref{thm:submod-marginal-is-submod},
for any $S \subseteq A_{\jhat}$, we get
$v_i(S \mid A_i) \le \sum_{c \in S} v_i(c \mid A_i) \le 0$.
This contradicts the fact that $\tau_{\jhat} > 0$.

\textbf{Case 2}: $v_i(\ghat \mid A_i) > 0$ for some $\ghat \in A_{\jhat}$.
\\ Then $\max(T) = \tau_{\jhat} \le v_i(\ghat \mid A_i)$.
By \ref{item:impl:propm-to-prop1:propm-goods}, we get
$w_iv_i([m]) < v_i(A_i) + \max(T) \le v_i(A_i \cup \{\ghat\})$.
But this contradicts \ref{item:impl:propm-to-prop1:unprop1-goods}.

Hence, it cannot happen that $i$ is PROPm-satisfied by $A$ but not PROP1-satisfied.
\end{proof}

\thmRecall{rthm:impl:prop-to-pessShare}
\begin{proof}
Let $P$ be agent $i$'s $\ell$-out-of-$d$-partition.
Assume \wLoG{} that $v_i(P_1) \le v_i(P_2) \le \ldots \le v_i(P_d)$.
Then by \cref{thm:prefix-sum-bound} and superadditivity of $v_i$, we get
\[ \loodM_i = \sum_{j=1}^{\ell} v_i(P_j)
    \le \frac{\ell}{d}\sum_{j=1}^d v_i(P_j) \le \frac{\ell}{d}v_i([m]). \]
Hence,
\[ \pessShare_i \defeq \sup_{1 \le \ell \le d: \ell/d \le w_i} \loodM_i
    \le \sup_{1 \le \ell \le d: \ell/d \le w_i} (\ell/d)v_i([m]) \le w_iv_i([m]).
    \qedhere \]
\end{proof}
